\documentclass[a5paper,10pt,headings=small]{scrartcl}
\usepackage[portrait,top=1.5cm,left=2cm,bottom=1.5cm,right=2cm]{geometry}
\usepackage[utf8]{inputenc}
\usepackage[T1]{fontenc}
\usepackage{textcomp} % für €
\usepackage{tabularx}
\usepackage{libertine} % kann man notfalls auch ignorieren, wenns nicht da ist
\usepackage{wasysym} % für \Square
\usepackage{multirow}
\usepackage{wrapfig}
\usepackage{graphicx}

\pagestyle{empty}
\setlength{\parindent}{0cm}
\setlength{\parskip}{1ex}
\newlength{\signskiplen}
\setlength{\signskiplen}{2ex}
\newcommand{\signskip}{\rule{0pt}{18pt}}
\newcommand{\hinweis}[1]{\emph{#1}}

\title{Mitgliedsantrag}

\begin{document}

% Header mit Logo und Anschrift
\newsavebox{\headerboxaddress}
\sbox{\headerboxaddress}{
  \parbox{.4\textwidth}{
    \small
    Vorstand des Stratum~0~e.\,V. \\
    Hamburger Straße 273a \\
    38112 Braunschweig \\
    vorstand@stratum0.org \\
    Fax: (0531) 287 69 246
  }
}
\newsavebox{\headerboxlogo}
\sbox{\headerboxlogo}{
  \parbox{.4\textwidth}{
    \hfill % damit rechtsbündig klappt
    \includegraphics[height=52pt]{stratum0-logo.eps}
  }
}
\begin{center}
\begin{tabular}{@{}p{0.4\textwidth}@{\phantom{m}}p{0.4\textwidth}}
  \multicolumn{1}{r}{
    \usebox{\headerboxlogo}
  }
  &
  \usebox{\headerboxaddress}
\end{tabular}
\end{center}

\vfill

% angepasster Titel... \maketitle ist Mist auf a5paper
\begin{center}
  \Large \titlefont \makeatletter \@title \makeatother
\end{center}

\emph{Mitgliedsanträge können auch formlos per E-Mail oder Fax gestellt werden,
unter Angabe der hier aufgeführten Daten.}

% Persönliche Daten
\section*{Persönliche Daten}
\begin{description}
  \item[Bürgerlicher Name:] \hrulefill% \signskip
  \item[E-Mail-Adresse:] \hrulefill \signskip
\end{description}
\hinweis{Alle Änderungen dieser Daten müssen dem Vorstand gemeldet werden. Falls
der Vorstand dich in wich\-ti\-gen Fällen nicht erreichen kann, hat dies unter
Umständen deinen Aus\-schluss als Mitglied zur Folge.}

% Mitgliedsbeitrag
\section*{Mitgliedsbeitrag}
\begin{description}
  \item[\Square] Ich möchte den vollen Beitrag von \textbf{20€ pro Monat} zahlen.
  \item[\Square] Ich gehe zur Schule, studiere, empfange Sozialgeld oder ALG~II
    einschließ\-lich Leistungen nach §~22 ohne Zuschläge oder nach §~24 SGB~II,
    oder beziehe Ausbildungsförderung nach dem BAföG und möchte den ermäßigten
    Beitrag von \textbf{12€ pro Monat} zahlen. Ein entsprechender Nachweis muss
    dem Vorstand auf Verlangen zugänglich gemacht werden.
  \item[\Square] Ich kann aus finanziellen Gründen den regulären Beitrag nicht
    aufbringen und beantrage nach §3~Abs.~3 der Beitragsordnung eine Ermäßigung
    des Betrags. Dieser Antrag gilt für ein Jahr und kann erneuert werden. \\
    \Square~Befreiung \hfill
    \Square~Ermäßigung auf: \hrulefill~€\hfill\phantom{a}
\end{description}

\vfill

Datum, Unterschrift: \hrulefill\signskip

\hinweis{Sobald der Vorstand über deinen Antrag entschieden hat, bekommst du per
E-Mail Nachricht über deine Aufnahme und erhältst weitere Informationen.}

\end{document}
% vim: set tw=80 et sw=2 ts=2:
